% 核心结论：串联电阻电感电容零输入响应的三种阻尼情形
% 图原型：纯几何／对比。物理网络按本题节点、端口和符号构建。
\documentclass[tikz,border=7pt]{standalone}
\usepackage{fontspec,xeCJK,amsmath,amssymb}
\IfFontExistsTF{Times New Roman}{\setmainfont{Times New Roman}}{\setmainfont{TeX Gyre Termes}}
\IfFontExistsTF{Heiti SC}{\setCJKmainfont[AutoFakeBold=1.5]{Heiti SC}}{\setCJKmainfont[AutoFakeBold=1.5]{FandolHei-Regular}}
\usetikzlibrary{arrows.meta,calc,angles,quotes}
\definecolor{ink}{HTML}{243745}
\definecolor{blue}{HTML}{275DAB}
\definecolor{green}{HTML}{177C63}
\definecolor{red}{HTML}{BD4935}
\definecolor{purple}{HTML}{715393}
\begin{document}
\begin{tikzpicture}[font=\fontsize{9.5}{12}\selectfont,text=ink,
 v/.style={line width=1.1pt,-{Stealth[length=2.4mm,width=1.5mm]}},
 axis/.style={ink!65,line width=.45pt,-{Stealth[length=1.7mm,width=1.1mm]}},
 aux/.style={ink!40,densely dashed,line width=.5pt},
 lab/.style={fill=white,inner sep=2pt},
 card/.style={draw=blue!25,fill=blue!3,rounded corners=3pt,text width=6cm,align=left,inner sep=10pt}]
\path[use as bounding box] (-8,-3.35) rectangle (8,3.65);
\node[font=\fontsize{10.5}{13}\selectfont\bfseries,text width=15cm,align=center] at (0.00000,3.22000) {M6-H14　串联电阻电感电容零输入响应的三种阻尼情形};
\begin{scope}[shift={(-4,0)}]
\draw[axis] (-2.50000,-0.65938) -- (2.72000,-0.65938);
\draw[axis] (-2.50000,-1.85000) -- (-2.50000,1.97000);
\node[] at (2.60000,-2.55000) {$\omega_0t$};
\node[] at (-2.25000,2.12000) {$q/q_0$};
\draw[blue,thick] plot[domain=0:10,samples=250] ({-2.5+0.51*(\x-0)},{-1.8+2.28125*((exp(-.3*\x)*(cos(deg(sqrt(.91)*\x))+.3/sqrt(.91)*sin(deg(sqrt(.91)*\x))))-(-0.5))});
\draw[green,thick] plot[domain=0:10,samples=250] ({-2.5+0.51*(\x-0)},{-1.8+2.28125*(((1+\x)*exp(-\x))-(-0.5))});
\draw[red,thick] plot[domain=0:10,samples=250] ({-2.5+0.51*(\x-0)},{-1.8+2.28125*((((2+sqrt(3))*exp((-2+sqrt(3))*\x)+(-2+sqrt(3))*exp((-2-sqrt(3))*\x))/(2*sqrt(3)))-(-0.5))});
\draw[ink] (-2.50000,-0.71938) -- (-2.50000,-0.59937);
\node[] at (-2.50000,-2.20000) {$0$};
\draw[ink] (0.05000,-0.71938) -- (0.05000,-0.59937);
\node[] at (0.05000,-2.20000) {$5$};
\draw[ink] (2.60000,-0.71938) -- (2.60000,-0.59937);
\node[] at (2.60000,-2.20000) {$10$};
\draw[ink] (-2.56000,-1.80000) -- (-2.44000,-1.80000);
\node[anchor=east] at (-2.72000,-1.80000) {$-0.5$};
\draw[ink] (-2.56000,-0.65938) -- (-2.44000,-0.65938);
\node[anchor=east] at (-2.72000,-0.65938) {$0$};
\draw[ink] (-2.56000,1.85000) -- (-2.44000,1.85000);
\node[anchor=east] at (-2.72000,1.85000) {$1.1$};
\node[text=blue] at (-1.70000,2.55000) {$\zeta=.3$};
\node[text=green] at (-0.30000,2.55000) {$\zeta=1$};
\node[text=red] at (1.10000,2.55000) {$\zeta=2$};
\end{scope}
\begin{scope}[shift={(4,0)}]
\node[card] at (0.00000,1.15000) {$\alpha=R/(2L)$};
\end{scope}
\begin{scope}[shift={(4,0)}]
\node[card] at (0.00000,-1.15000) {$\alpha\lessgtr\omega_0\text{ 决定根型}$};
\end{scope}
\node[text width=15cm,align=center] at (0.00000,-3.00000) {曲线用明确的归一化变量；零输入与零状态阶跃不可混淆。};

\end{tikzpicture}\end{document}